\addtocontents{toc}{\protect\Needspace{3\baselineskip}}
\originalchapter{8.821 课程详讲：近边界模式、稳定性与标度维数}
本章对应第 19 讲. 质量--维数关系不是把 $m$ 与 $\Delta$ 凭直觉相等, 而是近边界波动方程和源的尺度变换共同决定. 前置是曲背景 Klein--Gordon 方程、守恒内积和上一章的边界源. 先研究一个自由标量, 最后再比较矢量与应力; 背景取 $\mathrm{AdS}_{d+1}$ Poincaré 度规.

\originalsection{大 \texorpdfstring{$N$}{N} 为什么允许先解自由体场}
一般体作用量有 $1/(2\kappa^2)$ 的整体因子, $2\kappa^2=16\pi G_{d+1}$. 将场涨落相对于背景作适当正则重标度, 使二次项无 $1/\kappa^2$, 三次顶点便带 $\kappa$, 四次带 $\kappa^2$. 在以曲率为单位的 $\kappa\sim N^{-1}$ 极限, 有限幅度涨落首阶近似自由. 这说明先求模式与质量的合理性, 而不是在没有说明极限时任意忽略相互作用.

取二次 Lorentzian 作用量
\[
S=-\frac12\int\dd^{d+1}x\sqrt{-g}
(g^{MN}\partial_M\phi\partial_N\phi+m^2\phi^2),
\quad (\nabla^2-m^2)\phi=0.
\]
沿边界方向 Fourier 展开 $\phi(z,x)=\int\dd^dk\,\ee^{\ii k\cdot x}\phi(z,k)/(2\pi)^d$, $k^2=-\omega^2+\mathbf k^2$. 用 $\sqrt{-g}=(R/z)^{d+1}$, 得
\[
z^{d+1}\partial_z(z^{1-d}\partial_z\phi)
-k^2z^2\phi-m^2R^2\phi=0.
\]
近 $z=0$、固定有限 $k$ 时, $k^2z^2$ 比质量项与 Euler 导数项多两个 $z$ 次数, 可先解其主导幂律. 试 $\phi\sim z^\beta$, 得 $\beta(\beta-d)=m^2R^2$. 定义
\[
\nu=\sqrt{\frac{d^2}4+m^2R^2},\quad
\Delta_+=\frac d2+\nu,\quad \Delta_- =\frac d2-\nu,
\]
在两支非简并且尚未遇到有关 log 项时,
\[
\phi(z,x)=z^{\Delta_-}A(x)+z^{\Delta_+}B(x)+\cdots.
\]
省略项可含边界导数和更高幂; 当指数差为有关整数、或 $\nu=0$ 时还可能有 log, 不能把两支简单幂当作全部情形.

\originalsection{负质量平方在 AdS 不必意味着不稳定}
平直时空自由标量有 $\omega^2=\mathbf k^2+m^2$, 若 $m^2<0$, 足够长波令 $\omega^2<0$, 对应指数增长. AdS 的边界与曲率迫使模式支付梯度能量, 可抵消一定范围的负势能. 自由标量的 Breitenlohner--Freedman（BF）条件是
\[
m^2R^2\ge-\frac{d^2}4,
\]
它正是上面 $\nu$ 取实数的条件. 还须选适当的保能边界条件; BF 条件不保证任意混合边界作用量都没有新的不稳定.

可从径向方程直接看见临界值. 令 $\phi=z^{(d-1)/2}\psi(z)\ee^{-\ii\omega t+\ii\mathbf k\cdot\mathbf x}$, 方程化为
\[
-\psi''+\frac{\nu^2-1/4}{z^2}\psi
=(\omega^2-\mathbf k^2)\psi.
\]
对端点消失的光滑试探函数, 恒等式
\[
\int_0^\infty\left|\psi'-\frac\psi{2z}\right|^2\dd z
=\int_0^\infty|\psi'|^2\dd z
-\frac14\int_0^\infty\frac{|\psi|^2}{z^2}\dd z\ge0
\]
给 Hardy 下界. 因而 $\nu^2\ge0$ 时, 相应正规径向二次型至少在这类试探域非负; 更完整的自伴边界实现需采用合适的模式条件. 若 $\nu^2<0$, 逆平方吸引超过临界值, 可用靠近边界的函数压低梯度相对代价, 二次型失去下界, 对应 BF 不稳定. 这个推导解释了为什么只看 $m^2$ 的正负会错判 AdS 中的稳定性.

\originalsection{KG 内积的守恒和近边界收敛}
两个同质量解的流 $j^M=\ii(\phi_1^*\nabla^M\phi_2-\phi_2\nabla^M\phi_1^*)$ 散度为零, 因为两场方程中的质量项相消. 在空间型时片 $\Sigma$ 上积分该流, 得 KG 内积. 两时片内积的差由边界通量给出, 所以仅满足体方程还不够, 也需边界通量为零.

定 $t$ 片的测度乘单位法向时间分量按 $R^{d-1}z^{1-d}$ 缩放. 对有限频率模式 $\phi\sim z^\beta$, 近边界内积的径向收敛要求
\[
\int_0\dd z\,z^{2\beta+1-d}<\infty,
\qquad \beta>\frac{d-2}2.
\]
$\Delta_+$ 支对实 $\nu\ge0$ 满足; $\Delta_-$ 支要 $\nu<1$. 因而通常 $\nu\ge1$ 只有快衰减支可作有限 KG 范数的普通量子模式. $0<\nu<1$ 则两支都可有有限范数. $\nu=0$ 是重根需独立 log 处理, $\nu=1$ 慢支处于普通内积的对数边界, 不纳入本章简单的替代量子化窗口.

范数有限仍不自动给所有两解之间零通量. 展开两支后, 源/响应交叉项的通量正比于
\[
2\nu(A_1^*B_2-B_1^*A_2).
\]
选择 $A=0$ 的标准量子模式, 或窗口内 $B=0$ 的替代模式, 都使它消失. 合适 real/自伴的混合关系也可使通量为零, 但其能量稳定性还须另外检验. 因此“可归一化”是范数条件, “已选量子态空间”还包含允许的边界关系.

满足所选关系的一套正频正常模 $u_n$ 可构成场算符 $\phi=\sum_n(u_na_n+u_n^*a_n^\dagger)$. 这才是固定理论中建 Hilbert 空间的态. 另一支被固定为外部经典源时, 改变它是在改变背景耦合, 不是给同一无源 Hilbert 空间再多加一类不受约束粒子.

\originalsection{尺度权把源系数与算符维数联系起来}
在标准量子化中取 $\Delta=\Delta_+$, 源是
\[
\phi_0(x)=A(x)=\lim_{z\to0}z^{\Delta-d}\phi(z,x).
\]
它一般不是 $\phi(0,x)$ 的有限数值. 联合 AdS 等距缩放 $x'=\lambda x,z'=\lambda z$ 下, 标量满足 $\phi'(z',x')=\phi(z,x)$. 比较慢支系数,
\[
\phi'_0(x')=\lambda^{\Delta-d}\phi_0(x).
\]
边界源项 $\int\dd^dx\,\phi_0\mathcal O$ 保持同一形式, 所以 $\mathcal O'(x')=\lambda^{-\Delta}\mathcal O(x)$. 因此
\[
\boxed{m^2R^2=\Delta(\Delta-d),\qquad
\Delta=\frac d2+\sqrt{\frac{d^2}4+m^2R^2}}
\]
是标准量子化下的标量质量--维数关系. 快支系数 $B$ 也以 $\lambda^{-\Delta}$ 变换, 符合它与算符期望值的联系, 系数将在下一章从壳上作用量求出.

在 $0<\nu<1$ 的替代量子化中, 交换哪支定义源与状态, 算符维数为 $\Delta_-$. 此时源是快支 $B$ 的适当系数, 不是仍把 $A$ 称作同一个源. 两种量子化通常定义不同的边界理论, 不能只把同一场的一次读数随手改名而得到另一维数.

\begin{example}
在 $d=4$, $m^2R^2=-15/4$ 时 $\nu=1/2$, 标准维数 $\Delta_+=5/2$, 替代维数 $\Delta_-=3/2$, 都在简单双支窗口内. 对 $m^2R^2=-3$, $\nu=1$, 标准维数3, 慢支的普通 KG 范数处于对数边界, 不能直接当作前一个窗口例子. 对 $m^2R^2=-5$, 已低于 BF 的 $-4$, 问题不是选择两种实维数, 而是自由体标量本身不稳定.
\end{example}

\originalsection{相关、边缘与无关源的物理含义}
标准量子化下, $m^2=0$ 给 $\Delta=d$, 是边缘算符; BF 允许范围中的 $m^2<0$ 给 $\Delta<d$, 是相关算符; $m^2>0$ 给 $\Delta>d$, 是无关算符. 这里“相关”指源耦合在低能变重要的 RG 意义, 不是统计上的两变量相关系数.

相应慢支 $z^{d-\Delta}$ 在边界 $z\to0$ 分别常数、消失或增长. 对无关源, 有限强度源可能在紫外端改变原有渐近条件, 不能毫无回馈限制地使用固定 AdS 的小涨落近似. 本课程的关联函数计算取围绕零源的微扰泛函导数, 与这一限制相容.

\originalsection{矢量和应力的尺度不能照抄标量}
对质量 $m$ 的 transverse Proca 模, 径向方程为
\[
z^2A_\mu''+(3-d)zA_\mu'-(k^2z^2+m^2R^2)A_\mu=0.
\]
幂指数满足 $\beta(\beta-d+2)=m^2R^2$. 考虑一形式本身的坐标权后, 标准矢量算符维数为
\[
\Delta_V=\frac d2+\sqrt{\frac{(d-2)^2}4+m^2R^2},
\qquad m^2R^2=(\Delta_V-1)(\Delta_V-d+1).
\]
当 $m=0$, $A_\mu=a_\mu+b_\mu z^{d-2}+\cdots$, $\Delta_J=d-1$, 与守恒流的已知维数一致. Proca 的 longitudinal 与径向约束以及全部入选小问在 PS5.1 中给出, 这里先展示横向指数为何与标量不同.

应力源是 dimensionless 的边界度规变形, $\int h_{\mu\nu}T^{\mu\nu}$ 要有无量纲作用量, 得 $\Delta_T=d$. 同一原始体 metric 的 $z^{-2}$ 共形因子与 metric 涨落的 frame 需分别处理, 不能把标量的 $z^{d-\Delta}$ 直接移到所有 raw 张量分量上. 这完成了第 18 讲只提出源字典、尚未说明指数的那一步.
